\chapter{Proofs for litigation inference games}
\label{annex:p18-lis-burdens}

\begin{annexitem}{annex:p18.lis-grounded-game}{thm:p18-lis-grounded-game}{Grounded justification and the fixed-role game}
\begin{proof}
The defense operator \(F\) is monotone: a defender available in
\(S\) remains available in every \(T\supseteq S\).
Consequently the iteration from the empty set is increasing.
At a limit stage, any element already admitted remains admitted
at the next stage by monotonicity.
If the iteration never stabilized before an ordinal of cardinality
greater than \(|\mathcal A|\), choosing a newly admitted argument
at each strict increase would inject that ordinal into \(\mathcal A\).
Thus it stabilizes at a fixed point \(J\).
For every fixed point \(K\), transfinite induction gives
\(J_\alpha\subseteq K\); hence \(J\) is the least fixed point.
Every first-entry stage is a successor, because limits introduce
only unions.

First suppose \(A\in J\), with first-entry stage \(\alpha+1\).
For every \(BDA\), the defining property supplies a defender
\(C\in J_\alpha\) with \(CDB\).
Among all such defenders in \(J_\alpha\), choose one of minimum
entry rank. It strictly defeats \(B\).
Indeed, if \(BDC\), then the earlier-stage defense required for
\(C\)'s membership supplies \(C'\) defeating \(B\) at a smaller
entry rank, contradicting the choice of \(C\).
This proves the strict-defender lemma without assuming that
arbitrary defenders are strict.

At each opponent move, let proponent choose such a minimum-rank
strict defender of that move. Its rank is smaller than that of
the preceding proponent argument.
Use a fixed well-order to break ties if necessary.
The argument ranks at proponent turns therefore strictly decrease.
An ordinal admits no infinite strictly decreasing sequence, so
every compatible play is finite.
Proponent can answer every opponent move by the lemma and cannot
repeat a previous proponent argument, since repetition would repeat
its rank. The only possible terminal is consequently one in which
opponent cannot move. This is a winning strategy.

Conversely, suppose a proponent winning strategy exists.
Unfold its complete tree of legal histories, keeping all opponent
choices and the selected proponent response at each strategy node.
Every maximal branch is finite and ends at a proponent move.
The tree is well founded: an infinite descending path would give
an infinite play contrary to the definition of winning.
Work by well-founded induction on its proponent nodes.
At a leaf, the moved argument has no defeaters, because opponent
is allowed every defeating argument, including repetitions.
It belongs to \(F(\varnothing)\subseteq J\).

At any other proponent node with argument \(A'\), take an arbitrary
\(BDA'\). The opponent move \(B\) is available regardless of its
prior appearances. The winning strategy supplies a proponent
response \(C\), and the subtree below that response is winning.
The induction hypothesis gives \(C\in J\).
Its legal response strictly defeats \(B\), in particular defeats it.
Every defeater \(B\) of \(A'\) thus has a defender in \(J\),
so \(A'\in F(J)=J\).
This holds at the root and proves \(A\in J\).

The proof uses well-founded trees and ranks rather than a finite
depth bound. Under ordinary classical set theory with choice it
applies to arbitrary set-sized branching; no finite-branching
assumption was inserted.
\end{proof}
\end{annexitem}

\begin{annexitem}{annex:p18.lis-opponent-repetition}{thm:p18-lis-opponent-repetition}{Counterexample to forbidding opponent repeats}
\begin{proof}
Let the argument set be \(\{A,B,C,D,E\}\), with exactly the
following directed defeat edges:
\[
B\longrightarrow A,\quad C\longrightarrow B,\quad
D\longrightarrow C,\quad E\longrightarrow D,\quad
B\longrightarrow E .
\]
There are no reverse edges, so all five displayed defeats are strict.
Every argument has a defeater. Hence \(F(\varnothing)=\varnothing\),
and the grounded extension is empty.

If each player is forbidden to repeat its own arguments, the only
play starting at \(A\) is
\[
P:A,\ O:B,\ P:C,\ O:D,\ P:E .
\]
The only defeater of \(E\) is \(B\), but opponent has already used it.
Thus opponent has no legal move and proponent wins.
There were no alternative choices, so this is a winning strategy.
Nevertheless \(A\) is not grounded.

In the corrected game, opponent can repeat \(B\).
Proponent would then need \(C\), its only defender, but proponent
has already used \(C\); that repetition is forbidden.
The play therefore ends against proponent, in agreement with the
empty grounded extension.
\end{proof}
\end{annexitem}

\begin{annexitem}{annex:p18.lis-burden-domain}{thm:p18-lis-burden-domain}{Closure obstruction and literal profiles}
\begin{proof}
Let \(S,T\subseteq\mathcal L\) be closed under unrestricted
classical consequence. For any classical tautology \(\tau\),
both \(S\models_{\rm CL}\tau\) and \(T\models_{\rm CL}\tau\).
Closure gives \(\tau\in S\cap T\), contradicting the prohibition
on assigning the same formula to both parties.
This argument applies even if the intended generating sets of
burdens were empty. Consistency does not remove tautologies.
Thus unrestricted classical closure cannot satisfy the printed
separation requirement.

Now let \(Q\) contain only literals and let \(S\subseteq Q\)
be classically consistent. Extensivity gives
\(S\subseteq\operatorname{Cn}_Q(S)\).
If \(l\in Q\setminus S\), construct a valuation satisfying \(S\)
and falsifying \(l\).
If \(\overline l\in S\), every valuation satisfying \(S\)
already falsifies \(l\). If neither literal is in \(S\), assign
the underlying propositional atom the value that falsifies \(l\),
and satisfy every other literal of \(S\).
Consistency ensures these assignments do not conflict.
Therefore \(S\not\models_{\rm CL}l\).
It follows that \(\operatorname{Cn}_Q(S)=S\).

Choose any two consistent literal sets \(S_\pi,S_\delta\) such
that neither a literal nor its complement is assigned across
the two sets. Define \(B(h,p)=S_p\) at every history.
Consistency, relative closure and separation hold.
No burden disappears at any step, so both retention requirements
hold as well. This supplies nonempty profiles whenever, for
example, \(S_\pi=\{a\}\) and \(S_\delta=\{b\}\) concern different
atoms. Empty sets also give a profile.
These are witnesses for the corrected issue-relative domain,
not for the impossible unrestricted-classical condition.
\end{proof}
\end{annexitem}

\begin{annexitem}{annex:p18.lis-role-totality}{thm:p18-lis-role-totality}{Construction of historical subtree roles}
\begin{proof}
At the first move the role of every occurrence is specified as \(P\).
Suppose all roles at moves before \(i\) have been assigned uniquely.
For the root of \(A_i\), membership of its conclusion in the
current mover's burden set is a determinate input.
If membership holds, its role is \(P\).
Otherwise the root argument defeats \(A_{i-1}\) by the candidate
prerequisite. The unique target \(t(A_i,A_{i-1})\) has a uniquely
assigned historical role by the induction hypothesis.
Taking its opposite defines the root role.
The parent-inheritance clause is therefore never needed at the root.

Proceed through the finite tree in order of increasing depth.
At an occurrence \(U\), the explicit-burden test is determinate.
If it fails, whether its subtree defeats the preceding argument
is also an input fact. If it does, the unique shared target gives
the required opposite role. If it does not, \(U\) is a proper
occurrence, so it has exactly one immediate parent, whose role
was already assigned at smaller depth.
The ordered cases are mutually exclusive as operative cases
and exhaust all occurrences. Induction on depth and then on
move number proves existence and uniqueness.

The occurrence convention matters: identical subderivations can
appear twice under different parents, and their inherited roles
need not be identified merely because their textual contents match.

For the target restriction, start with an argument \(Z\) by
\(\pi\), whose occurrences all have role \(P\).
Let \(\delta\)'s next argument have root conclusion \(q\) and
incomparable subarguments \(U,V\) with conclusions \(u,v\).
Assign \(\{q,u\}\) to \(\delta\) and no issues to \(\pi\);
let \(V\) strictly defeat \(Z\), with this defeat propagated
to the whole moved argument and no additional target in \(Z\).
These are consistent literal burdens and a legal \(P\)-role
second move. The root and \(U\) have role \(P\) from their
explicit burdens. \(V\) has role \(O\) by attacking the earlier
\(P\)-role target without an explicit burden.
Now permit a subsequent argument with distinct branches concluding
\(\neg u\) and \(\neg v\), attacking both \(U\) and \(V\).
If it has no explicit burden, choosing \(U\) as its target would
give role \(O\), while choosing \(V\) would give role \(P\).
There is no unique smallest target and no unique role from the
source rule. The declared single-target domain excludes precisely
this case rather than selecting an arbitrary target.
\end{proof}
\end{annexitem}

\begin{annexitem}{annex:p18.lis-winning-construction}{thm:p18-lis-winning-construction}{Winning sets for the full history arena}
\begin{proof}
The operator \(\Phi_p\) is monotone because enlarging \(X\)
preserves both the existential and universal successor conditions.
Iteration from the empty set is increasing, with union at limits.
It stabilizes by the same set-cardinality argument as for the
grounded operator, now on the set \(H\) of finite histories.
Call its least fixed point \(W_p\).
Each first-entry rank is a successor rank.

For \(h\in W_p\), select a strategy as follows.
If \(h\) is a terminal in \(T_p\), the win has already occurred.
Otherwise let its first-entry rank be \(\alpha+1\).
At a \(p\)-turn, the existential clause supplies a successor
in \(X_\alpha\); choose one, breaking ties by a fixed well-order.
At a \(\bar p\)-turn, the universal clause ensures that every
successor is in \(X_\alpha\), and that a successor exists.
Thus every move before termination strictly decreases the rank.
There cannot be an infinite compatible play. The only terminal
histories admitted to \(W_p\) are in \(T_p\): a terminal with
turn \(p\) satisfies neither existential nor guarded universal
clause and is not in \(T_p\).
Every resulting terminal is consequently a win for \(p\).

For the converse, let \(h\notin W_p\).
If it is a \(p\)-turn, no successor lies in \(W_p\), since otherwise
the existential fixed-point clause would place \(h\) in \(W_p\).
If it is a \(\bar p\)-turn, it cannot be terminal:
such a terminal would be in \(T_p\).
Moreover, some successor lies outside \(W_p\), because otherwise
the nonempty universal clause would place \(h\) in \(W_p\).
Let \(\bar p\) choose such an outside successor at its turns.
Against every strategy of \(p\), the play stays outside \(W_p\).
If it ends, it cannot end in \(T_p\).
If it is infinite, it is not a terminal win for \(p\).
Hence \(p\) has no strategy forcing the required win.
This proves the equivalence on the entire history-dependent arena.

If \(h\) belonged to both winning sets, combine the two supposed
winning strategies. They determine a compatible maximal play.
Each strategy requires that play to terminate with the other
party unable to move. A finite terminal cannot have both parties
as the next player, and an infinite play satisfies neither winning
condition. The supposition is impossible.

Finally, if the supplied arena has \(N\) states, each strict
increase adds at least one previously absent state. There are
at most \(N\) such increases. Once an iteration makes no change,
the fixed point has been reached. No inference from finite
\(\mathcal A\) alone to finite \(H\) is used.
\end{proof}
\end{annexitem}

\begin{annexitem}{annex:p18.lis-grounded-separation}{thm:p18-lis-grounded-separation}{Litigation winning is not grounded membership}
\begin{proof}
Use the three arguments and constant profile in the first worked
game: \(BDA\), \(BDC\), \(CDB\), with no other pair;
\(\delta\) owns \(b\), and \(\pi\) has no explicit burden.
The conclusions are \(\neg b,b,\neg b\), respectively.
The profile is admissible by the literal-profile theorem.
All trees have only a root, so targets are unique.

After initial \(A\), the only possible response is \(B\).
Its explicit burden gives role \(P\), and it strictly defeats \(A\).
The only reply defeating \(B\) is \(C\).
Its conclusion is not in \(\pi\)'s burden set; the previous
target has role \(P\), so \(C\) has role \(O\).
Its mutual defeat is legal.
The only defeater of \(C\) is \(B\), whose reuse by \(\delta\)
is prohibited because its earlier role was \(P\).
Thus the unique maximal litigation play ends with \(\pi\) winning.

Every argument in this defeat graph has a defeater.
Therefore \(F(\varnothing)=\varnothing\), its grounded extension
is empty, and in particular \(A\notin J\).
This proves separation without relying on an intuitive legal
assessment or equating a merely selected branch with a strategy.
\end{proof}
\end{annexitem}

\begin{annexitem}{annex:p18.lis-protected-production}{thm:p18-lis-protected-production}{Protected rejection of an evidential argument}
\begin{proof}
Since \(U\) has no defeaters, \(U\in F(\varnothing)\subseteq J\).
Because \(UDE\), membership of \(E\) in \(F(J)\) would require
some member of \(J\) to defeat \(U\), which is impossible.
Thus \(E\notin J\), and it is defeated by the justified \(U\);
it is overruled.

If a reply with \(U\) is legal at a given history, the next party
must defeat \(U\) to continue, whether its role would be \(P\)
or \(O\). No such argument exists. The play is terminal and
the mover of \(U\) wins that branch.
The qualification that the move is legal includes the applicable
history restriction; the theorem does not authorize an otherwise
forbidden repetition.
\end{proof}
\end{annexitem}
